% kap10.tex — Hamiltonsk feltteori
% Hamilton-Jacobi-kompendiet

\chapter{Hamiltonsk feltteori}
\label{kap:felter-hamilton}

Vi slutter kompendiet, hvor vi startede: med Legendre-
transformationen (kapitel~\ref{kap:hamilton}), nu anvendt på felter i
stedet for partikler. Dette fuldender den klassiske del af broen mod
kvantefeltteori, og er selve kompendiets sidste skridt.

\section{Kanonisk impulstæthed}

I fuld analogi med \eqref{eq:def-p} definerer vi den \emph{kanoniske
impulstæthed}\index{kanonisk impulstæthed}, konjugeret til feltet $\varphi(x,t)$, som
\begin{equation}
  \boxed{\pi(x,t) \equiv \pdv{\mathcal{L}}{\dot\varphi}.}
  \label{eq:def-pi}
\end{equation}
Dette er den kontinuerte analog til $p_i=\partial L/\partial\dot q_i$
— hvor det diskrete indeks $i$ er blevet til det kontinuerte $x$.
Ligesom i det diskrete tilfælde kræver Legendre-transformationen, at
\eqref{eq:def-pi} kan inverteres, dvs.\ løses for $\dot\varphi$ som
funktion af $(\varphi,\pi,\varphi')$ — hvilket er opfyldt for alle
Lagrange-tætheder i dette kompendium, hvor $\mathcal{L}$ er kvadratisk
i $\dot\varphi$.

\section{Hamilton-tæthed og Hamiltons ligninger for felter}

Vi definerer \emph{Hamilton-tætheden}\index{Hamilton-tæthed} ved Legendre-transformation,
netop som i \eqref{eq:def-H}:
\begin{equation}
  \boxed{\mathcal{H}(\varphi,\pi,\varphi') \equiv \pi\dot\varphi - \mathcal{L},}
  \label{eq:def-H-felt}
\end{equation}
hvor $\dot\varphi$ på højresiden udtrykkes ved $(\varphi,\pi,\varphi')$
via inversionen af \eqref{eq:def-pi}. Bemærk, at dette er netop den
samme størrelse, vi allerede kaldte $\mathcal{H}$ i
\eqref{eq:kontinuitetsligning} i forrige kapitel — vi giver den nu en
systematisk definition som funktion af $(\varphi,\pi,\varphi')$ i
stedet for $(\varphi,\dot\varphi,\varphi')$, ganske som $H(q,p,t)$
afløste $H(q,\dot q,t)$ i kapitel~\ref{kap:hamilton}. Den samlede
Hamilton-funktion er $H=\int\mathcal{H}\,dx$.

\subsection{Udledning af bevægelsesligningerne}

Vi udleder Hamiltons ligninger for felter via det modificerede
Hamiltons princip (afsnit~\ref{sec:hamiltons-princip} i
kapitel~\ref{kap:kanoniske-transf}, nu udvidet til felter):
\begin{equation}
  S = \int dt\int dx\,\bigl[\pi\dot\varphi - \mathcal{H}(\varphi,\pi,\varphi')\bigr],
\end{equation}
med uafhængige, arbitrære variationer $\delta\varphi(x,t)$ og
$\delta\pi(x,t)$ (begge nul ved $t_1,t_2$ for $\delta\varphi$'s
tidsrand; ingen tidsrandbetingelse er nødvendig for $\delta\pi$, jf.\
opgave 1 i kapitel~\ref{kap:kanoniske-transf}). Variationen giver
\begin{equation}
  \delta S = \int dt\int dx\left[\pi\,\delta\dot\varphi + \dot\varphi\,\delta\pi
    - \pdv{\mathcal{H}}{\varphi}\delta\varphi - \pdv{\mathcal{H}}{\pi}\delta\pi
    - \pdv{\mathcal{H}}{\varphi'}\delta\varphi'\right].
\end{equation}
Vi integrerer $\pi\,\delta\dot\varphi$ partielt i $t$ (randtermen
forsvinder, som altid) og $-(\partial\mathcal{H}/\partial\varphi')\delta\varphi'$
partielt i $x$ (randtermen forsvinder ved strengens endepunkter, eller
i det uendelige):
\begin{equation}
  \delta S = \int dt\int dx\left\{\left[-\dot\pi - \pdv{\mathcal{H}}{\varphi}
    + \dv{}{x}\pdv{\mathcal{H}}{\varphi'}\right]\delta\varphi
    + \left[\dot\varphi - \pdv{\mathcal{H}}{\pi}\right]\delta\pi\right\}.
\end{equation}
Da $\delta\varphi$ og $\delta\pi$ er uafhængige og arbitrære, må hver
koefficient forsvinde for sig. Dette giver \emph{Hamiltons ligninger
for felter}\index{Hamiltons ligninger!for felter}:
\begin{equation}
  \boxed{\dot\varphi = \pdv{\mathcal{H}}{\pi}, \qquad
  \dot\pi = -\pdv{\mathcal{H}}{\varphi} + \dv{}{x}\pdv{\mathcal{H}}{\varphi'}.}
  \label{eq:hamilton-ligninger-felt}
\end{equation}
Den anden ligning har et ekstra led, $d/dx(\partial\mathcal{H}/\partial\varphi')$,
i forhold til den diskrete form \eqref{eq:hamiltons-ligninger} — en
direkte konsekvens af, at $\mathcal{H}$ (i modsætning til $H$ for en
partikel) også kan afhænge af feltets \emph{rumlige} afledte
$\varphi'$.

\section{Eksempel: strengen i Hamilton-billede}

\begin{eksempel}
For strengen er $\pi=\partial\mathcal{L}/\partial\dot\varphi=\mu\dot\varphi$,
så $\dot\varphi=\pi/\mu$, og
\begin{equation}
  \mathcal{H} = \pi\dot\varphi - \mathcal{L} = \pi\cdot\frac{\pi}{\mu}
    - \left(\frac12\mu\left(\frac{\pi}{\mu}\right)^2 - \frac12\tau(\varphi')^2\right)
  = \frac{\pi^2}{2\mu} + \frac12\tau(\varphi')^2,
\end{equation}
netop kinetisk plus potentiel energitæthed, som i
\eqref{eq:kontinuitetsligning}, nu udtrykt ved $\pi$ i stedet for
$\dot\varphi$. Hamiltons ligninger \eqref{eq:hamilton-ligninger-felt}
giver
\begin{equation}
  \dot\varphi = \pdv{\mathcal{H}}{\pi} = \frac{\pi}{\mu}, \qquad
  \dot\pi = -\pdv{\mathcal{H}}{\varphi} + \dv{}{x}\pdv{\mathcal{H}}{\varphi'}
          = 0 + \dv{}{x}(\tau\varphi') = \tau\varphi''.
\end{equation}
Den første er blot definitionen af $\pi$ igen; kombineres de to
($\dot\pi=\mu\ddot\varphi=\tau\varphi''$), genfindes bølgeligningen
\eqref{eq:boelgeligning}, $\ddot\varphi=v^2\varphi''$ — som forventet,
og som et konsistenstjek af hele Legendre-transformationen for feltet.
\end{eksempel}

\section{Eksempel: det komplekse felt og ladningen som generator}
\label{sec:komplekst-felt-hamilton}

\begin{eksempel}
Strengen er et reelt felt; for at se, hvordan Hamilton-billedet
håndterer et felt med \emph{intern} struktur, genbesøger vi det
komplekse felt $\psi(x,t)$ fra
afsnit~\ref{sec:ladning} i kapitel~\ref{kap:felter}, med
Lagrange-tæthed \eqref{eq:L-komplekst-felt},
$\mathcal{L}=\mu\dot\psi^\ast\dot\psi-\tau\psi^{\ast\prime}\psi'$.
Da $\psi$ og $\psi^\ast$ behandles som uafhængige felter, får de hver
sin kanoniske impulstæthed:
\begin{equation}
  \pi \equiv \pdv{\mathcal{L}}{\dot\psi} = \mu\dot\psi^\ast, \qquad
  \pi^\ast \equiv \pdv{\mathcal{L}}{\dot\psi^\ast} = \mu\dot\psi,
\end{equation}
(bemærk krydskoblingen: $\pi$, den til $\psi$ konjugerede impuls, er
proportional med $\dot\psi^\ast$, \emph{ikke} $\dot\psi$ — en
konsekvens af, at $\mathcal{L}$ blander $\psi$ og $\psi^\ast$).
Inverteres disse, $\dot\psi=\pi^\ast/\mu$, $\dot\psi^\ast=\pi/\mu$, og
Hamilton-tætheden bliver
\begin{equation}
  \mathcal{H} = \pi\dot\psi + \pi^\ast\dot\psi^\ast - \mathcal{L}
  = \frac{\pi\pi^\ast}{\mu} + \tau\,\psi^{\ast\prime}\psi'.
\end{equation}
Hamiltons ligninger \eqref{eq:hamilton-ligninger-felt} (nu ét par for
hvert af de to felter) reproducerer, ganske som for strengen, hver sin
bølgeligning for $\psi$ og $\psi^\ast$ — et konsistenstjek, vi
overlader til opgave~1.

\textbf{Ladningen som Poisson-generator.} Genkald Noether-ladningstætheden
$\rho=i(\psi\,\partial\mathcal{L}/\partial\dot\psi-\psi^\ast\,\partial\mathcal{L}/\partial\dot\psi^\ast)$
fra \eqref{eq:ladnings-kontinuitet}. Udtrykt ved de kanoniske impulser
er dette netop $\rho=i(\psi\pi-\psi^\ast\pi^\ast)$, så den totale
bevarede ladning er
\begin{equation}
  Q \equiv \int\rho\,dx = i\int\bigl(\psi\pi-\psi^\ast\pi^\ast\bigr)\,dx.
  \label{eq:Q-kanonisk}
\end{equation}
\index{Noether-ladning som generator}Vi beregner nu $\{\psi(x),Q\}$ med de fundamentale felt-Poisson-
parenteser \eqref{eq:felt-poisson} (udledt i næste afsnit, men brugt
her et skridt forud, da eksemplet naturligt hører sammen med $Q$):
kun leddet med $\{\psi(x),\pi(y)\}=\delta(x-y)$ i integranden bidrager,
\begin{equation}
  \{\psi(x),Q\} = i\int\psi(y)\,\{\psi(x),\pi(y)\}\,dy
  = i\int\psi(y)\,\delta(x-y)\,dy = i\,\psi(x).
  \label{eq:Q-genererer}
\end{equation}
Sammenlign nu med $U(1)$-transformationens infinitesimale form fra
kapitel~\ref{kap:felter}, $\delta\psi=i\varepsilon\psi$: resultatet
\eqref{eq:Q-genererer} viser præcis, at
\begin{equation}
  \boxed{\delta\psi = \varepsilon\,\{\psi,Q\},}
\end{equation}
altså den \emph{samme} formel, $\delta f=\varepsilon\{f,G\}$, som vi
udledte for infinitesimale kanoniske transformationer i
kapitel~\ref{kap:kanoniske-transf} — nu med Noether-ladningen $Q$
selv i rollen som generator $G$. Dette er ikke en tilfældighed, men
selve den dybere sandhed bag Noethers sætning i det Hamiltonske
billede: \emph{den bevarede størrelse, symmetrien giver anledning
til, er samtidig, via Poisson-parentesen, generatoren af symmetrien
selv} — ganske som $H$ selv genererer tidsudvikling
(kapitel~\ref{kap:kanoniske-transf}). Under den kanoniske
kvantiseringsregel, vi indfører i næste afsnit,
$\{\cdot,\cdot\}\to\frac{1}{i\hbar}[\cdot,\cdot]$, bliver $Q$ til
ladningsoperatoren $\hat Q$, og \eqref{eq:Q-genererer} bliver til
kommutatoren $[\hat\psi,\hat Q]=-\hbar\,\hat\psi$: $\hat Q$ genererer
$U(1)$-fasetransformationen på feltoperatoren $\hat\psi$, ganske som
$\hat H$ genererer tidsudvikling. Dette er en af udgangspunkterne
for det senere kvantefeltteori-kompendiums behandling af ladede
felter.
\end{eksempel}

\section{Poisson-parenteser for felter}

Vi generaliserer til sidst Poisson-parentesen (kapitel~\ref{kap:poisson})
til felter. Husk fra den diskrete kæde i
kapitel~\ref{kap:felter}, at $L=\sum_i a\,\mathcal{L}_i$, så den
diskrete kanoniske impuls er $p_i=\partial L/\partial\dot\varphi_i =
a\,\partial\mathcal{L}/\partial\dot\varphi_i = a\,\pi(x_i)$. De
fundamentale parenteser \eqref{eq:fundamentale-parenteser},
$\{\varphi_i,p_j\}=\delta_{ij}$, bliver derfor
$\{\varphi(x_i),\pi(x_j)\} = \delta_{ij}/a$. I kontinuumsgrænsen
$a\to0$ er $\delta_{ij}/a$ netop den diskrete repræsentation af
Diracs deltafunktion (arealet under $\delta_{ij}/a$, summeret med
vægt $a$ over $j$, giver $1$ — nøjagtig deltafunktionens
definerende egenskab). Vi finder derfor de \emph{fundamentale
felt-Poisson-parenteser}\index{Poisson-parentes!felt-}
\begin{equation}
  \boxed{\{\varphi(x),\pi(y)\} = \delta(x-y), \qquad
  \{\varphi(x),\varphi(y)\} = 0, \qquad \{\pi(x),\pi(y)\} = 0,}
  \label{eq:felt-poisson}
\end{equation}
evalueret ved samme tidspunkt $t$ for begge argumenter (såkaldte
\emph{ligetids-parenteser}).

\begin{bemaerkning}
Ligesom de fundamentale parenteser $\{q_i,p_j\}=\delta_{ij}$ for en
partikel er den klassiske forløber for kommutationsrelationen
$[\hat q_i,\hat p_j]=i\hbar\delta_{ij}$ (kapitel~\ref{kap:poisson}), er
\eqref{eq:felt-poisson} den klassiske forløber for kvantefeltteoriens
fundamentale ligetids-kommutationsrelation
\begin{equation}
  \bigl[\hat\varphi(x),\hat\pi(y)\bigr] = i\hbar\,\delta(x-y).
\end{equation}
Dette er selve udgangspunktet for kanonisk feltkvantisering: reglen
er, ganske som for et system med endeligt mange frihedsgrader, at
erstatte den klassiske Poisson-parentes med kommutatoren via
\begin{equation}
  \{f,g\} \;\longrightarrow\; \frac{1}{i\hbar}\bigl[\hat f,\hat g\bigr]
  \qquad\text{(kanonisk kvantisering}\index{kanonisk kvantisering}\text{)},
\end{equation}
anvendt her på selve feltet og dets impulstæthed. Et klassisk felt
$\varphi(x,t)$ (en almindelig, talværdi funktion) bliver herved til en
\emph{feltoperator} $\hat\varphi(x,t)$ — det centrale, nye objekt i
kvantefeltteorien, uden modstykke i den kvantemekanik, vi kender fra
[[kvantemekanik-kompendiet]] (hvor $\hat q_i$ er en operator, men det
klassiske udgangspunkt $q_i(t)$ er ét enkelt tal for hver frihedsgrad,
ikke en hel funktion af rummet). Vi udfører ikke selve kvantiseringen
her; det er det senere QFT-kompendiums opgave, og eksemplet ovenfor
med ladningen $Q$ som generator viser allerede, hvordan symmetri og
bevarelseslov (kapitel~\ref{kap:felter}) sammen med den kanoniske
struktur (dette kapitel) tilsammen forbereder netop det skridt. Men
med \eqref{eq:felt-poisson} har vi nået præcis det punkt, hvorfra
kvantiseringen tager sit udspring — og dermed også det punkt, hvor
dette kompendiums rejse ender.
\end{bemaerkning}

\begin{figure}[htbp]
  \centering
  \begin{tikzpicture}[scale=1.0]
    % Diskret faserum-gitter (q_i,p_i)
    \foreach \x/\lbl in {0/1,1.3/2,2.6/3,3.9/4} {
      \draw[-{Latex}] (\x,0) -- (\x,1.3);
      \filldraw[fill=blue!25,draw=blue!60!black] (\x,0.8) circle (0.08);
      \node[below] at (\x,-0.15) {\footnotesize$(\varphi_\lbl,\pi_\lbl)$};
    }
    \node at (1.95,1.8) {diskret: par $(\varphi_i,\pi_i)$ for hvert $i$};
    % Pil ned til kontinuert felt-par
    \draw[-{Latex},thick] (1.95,-0.6) -- (1.95,-1.1);
    \node at (3.1,-0.85) {$a\to0$};
    % Kontinuert par af felter
    \draw[thick,green!55!black] plot[smooth] coordinates
      {(0,-2.0) (1.0,-1.7) (2.0,-2.1) (3.0,-1.75) (3.9,-2.0)};
    \node[green!55!black] at (1.95,-2.6) {kontinuert: felt-par $(\varphi(x),\pi(x))$};
  \end{tikzpicture}
  \caption{Overgangen fra et diskret sæt kanonisk konjugerede par
  $(\varphi_i,\pi_i)$, ét for hver masse i kæden, til et kontinuert
  felt-par $(\varphi(x),\pi(x))$ — den kanoniske analog til
  kontinuumsgrænsen i figur~\ref{fig:kaede-til-streng}.}
  \label{fig:felt-poisson}
\end{figure}

\section{Opsummering}

\begin{itemize}
  \item Den kanoniske impulstæthed er $\pi=\partial\mathcal{L}/\partial\dot\varphi$;
        Hamilton-tætheden er $\mathcal{H}=\pi\dot\varphi-\mathcal{L}$.
  \item Hamiltons ligninger for felter er
        $\dot\varphi=\partial\mathcal{H}/\partial\pi$,
        $\dot\pi=-\partial\mathcal{H}/\partial\varphi + d/dx(\partial\mathcal{H}/\partial\varphi')$
        — sidste led skyldes $\mathcal{H}$'s afhængighed af den
        rumlige afledte $\varphi'$.
  \item For strengen: $\mathcal{H}=\pi^2/(2\mu)+\frac12\tau(\varphi')^2$,
        og Hamiltons ligninger reproducerer bølgeligningen.
  \item For det komplekse $U(1)$-felt: Noether-ladningen
        $Q=i\int(\psi\pi-\psi^\ast\pi^\ast)dx$ genererer selv symmetrien
        via Poisson-parentesen, $\delta\psi=\varepsilon\{\psi,Q\}$ —
        samme mønster som $H$ genererer tidsudvikling.
  \item De fundamentale felt-Poisson-parenteser er
        $\{\varphi(x),\pi(y)\}=\delta(x-y)$ — den klassiske forløber
        for kvantefeltteoriens ligetids-kommutator
        $[\hat\varphi(x),\hat\pi(y)]=i\hbar\delta(x-y)$, via
        kvantiseringsreglen $\{\cdot,\cdot\}\to\frac{1}{i\hbar}[\cdot,\cdot]$.
\end{itemize}

\begin{opgave}
Udled, ved at bruge Hamiltons ligninger \eqref{eq:hamilton-ligninger-felt}
på $\mathcal{H}=\pi\pi^\ast/\mu+\tau\psi^{\ast\prime}\psi'$
(afsnit~\ref{sec:komplekst-felt-hamilton}) for hvert af de to par
$(\psi,\pi)$ og $(\psi^\ast,\pi^\ast)$, at både $\psi$ og $\psi^\ast$
opfylder bølgeligningen $\ddot\psi=v^2\psi''$ (med $v^2=\tau/\mu$),
som forventet fra kapitel~\ref{kap:felter}.
\end{opgave}

\begin{opgave}
Genudfør udregningen \eqref{eq:Q-genererer} for $\psi^\ast$ i stedet
for $\psi$: vis, at $\{\psi^\ast(x),Q\}=-i\psi^\ast(x)$, og bekræft,
at dette stemmer med $\delta\psi^\ast=-i\varepsilon\psi^\ast$
(kapitel~\ref{kap:felter}) via $\delta\psi^\ast=\varepsilon\{\psi^\ast,Q\}$.
\end{opgave}

\section*{Kompendiets konklusion}
\addcontentsline{toc}{section}{Kompendiets konklusion}

Vi har nu tilbagelagt hele den rejse, kapitelstrukturen lagde ud: fra
Newtons love og D'Alemberts princip (kapitel~\ref{kap:newton-variation}),
over Lagrange- og Hamilton-formalismen
(kapitel~\ref{kap:lagrange}--\ref{kap:hamilton}), Poisson-parenteser
og kanoniske transformationer
(kapitel~\ref{kap:poisson}--\ref{kap:kanoniske-transf}), til
Hamilton-Jacobi-teorien og virknings-vinkel-variablene
(kapitel~\ref{kap:hj}--\ref{kap:virkning-vinkel}) — som i
kapitel~\ref{kap:adiabatisk} førte os direkte til Bohr-Sommerfeld-
kvantiseringen, det tankegods, der reelt var på plads, da
kvantemekanikken blev født. De to sidste kapitler
(\ref{kap:felter}--\ref{kap:felter-hamilton}) udvidede hele
formalismen fra endelige til uendelige frihedsgrader — klassiske
felter — og nåede frem til feltenes fundamentale Poisson-parenteser,
det naturlige udspringspunkt for kanonisk feltkvantisering.

Kompendiets erklærede mål — at nå frem til den ældre kvanteteoris
begrebsapparat, uden selv at udføre kvantiseringen — er hermed
opnået. To retninger står åbne herfra, begge allerede antydet
undervejs:

\begin{itemize}
  \item Den semiklassiske grænse ($\psi\sim e^{iS/\hbar}$, WKB,
        Bohr-Sommerfeld) forbinder til den fulde, korrekte
        kvantemekanik, som er [[kvantemekanik-kompendiet]]s emne.
  \item Kanonisk feltkvantisering ($\varphi\to\hat\varphi$,
        $\{\cdot,\cdot\}\to\frac{1}{i\hbar}[\cdot,\cdot]$, anvendt på
        felternes Poisson-algebra \eqref{eq:felt-poisson})
        fører til kvantefeltteori og videre til spredningsteori — det
        senere, selvstændige kompendiums emne.
\end{itemize}
